\section*{Выводы}
\addcontentsline{toc}{section}{Выводы}

В ходе выполнения курсовой работы была исследована система управления модулем поворота промышленного робота с ПИ-регулятором. На основании математической модели механической и электрической частей системы было получено характеристическое уравнение шестого порядка, зависящее от коэффициентов регулятора \(k_{\text{п}}\) и \(k_{\text{и}}\).

С использованием критериев Стодолы и Рауса--Гурвица была построена область устойчивости системы в пространстве параметров регулятора. Дополнительно для устойчивых точек были учтены технические ограничения на максимальный крутящий момент и управляющее воздействие:
\[
|M_{\text{м}}|_{\max} \le 150 \text{ Нм},
\qquad
|u|_{\max} \le 220 \text{ В}.
\]

В результате была получена допустимая область параметров регулятора, удовлетворяющая всем ограничениям. Проведённый анализ показал, что определяющим ограничением является максимальный динамический момент, тогда как ограничение по управляющему сигналу выполняется с запасом.

Для допустимой области были рассчитаны длительности переходных процессов и построены линии равных длительностей. Минимальная длительность переходного процесса на расчётной сетке составила
\[
t_{\text{п,min}} = 1.1 \text{ с}
\]
при параметрах
\[
k_{\text{п}} = 0.51,
\qquad
k_{\text{и}} = 1.01.
\]

Для выбора оптимальной настройки регулятора были исследованы интегральные показатели качества:
\[
J_{\psi}
=
\int_0^{t_{\text{п}}}
\psi_{\text{м}}^2(t)\,dt,
\qquad
J_u
=
\int_0^{t_{\text{п}}}
u^2(t)\,dt,
\]
а также совмещённый функционал
\[
J_{\text{оп}}
=
J_{\psi}
+
\rho J_u.
\]

Были определены характерные точки оптимальной кривой:
\[
\min J_{\psi}:
\qquad
k_{\text{п}} = 1.01,
\qquad
k_{\text{и}} = 2.01,
\]

\[
\min J_{\text{оп}}:
\qquad
k_{\text{п}} = 6.01,
\qquad
k_{\text{и}} = 3.01,
\]

\[
\min J_u:
\qquad
k_{\text{п}} = 12.4,
\qquad
k_{\text{и}} = 0.01.
\]

В результате численного перебора весового множителя было получено значение
\[
\rho = 0.24843,
\]
при котором минимум совмещённого функционала достигается в точке
\[
k_{\text{п}} = 6.01,
\qquad
k_{\text{и}} = 3.01.
\]

Для выбранной настройки были получены следующие показатели качества:
\[
t_{\text{п}} = 4.6 \text{ с},
\qquad
J_{\psi} = 10.0383,
\qquad
J_u = 32.8618,
\qquad
J_{\text{оп}} = 18.2021.
\]

Анализ переходных процессов показал, что данная настройка обеспечивает устойчивое движение системы, допустимый уровень динамического момента и умеренную величину управляющего сигнала при приемлемом быстродействии.

Таким образом, выбранная точка минимума функционала \(J_{\text{оп}}\) обеспечивает компромисс между точностью регулирования, быстродействием системы и величиной управляющего воздействия и может рассматриваться как оптимальная настройка регулятора для исследуемой системы.
